\scanpage{120}
% ms. line 01
\uline{Sur une certaine fonction holomorphe \deleted{\unsure{de variable|définie}} \inserted{sur l'}espace de Banach $l^1$.}

% ms. line 02
Soit, pour $\lambda=(\lambda_i)\in l^1$ et tout entier $n>0$
% ms. line 03
\[ \alpha_n(\lambda)=\sum_{i_1<\ldots<i_n}\lambda_{i_1}\ldots\lambda_{i_n} \]
% ms. line 03a
On a
% ms. line 04
\[ \alpha_n(\lambda)\leq\frac{\|\lambda\|_1^n}{n!} \]
% ms. line 05
les $\alpha_n(\lambda)$ sont les coefficients de la fonction entière
% ms. line 06
de genre zéro
% ms. line 07
\[ f(\unsure{\lambda|n|z})=\prod_{\unsure{1|i}}(1+\lambda_i z) \]
% ms. line 08
et on obtient, à un facteur constant près, les
% ms. line 09
coefficients les plus généraux d'une fonction entière
% ms. line 10
de genre zéro.

% ms. line 11
\uline{Proposition} Pour tout $\lambda\in l^1$, la série $\sum_0^\infty n!\alpha_n(\lambda)$
% Vertical rule beside the proposition.
% ms. line 12
est absolument convergente, soit $S(\lambda)$ sa somme.
% ms. line 13
La fonction $\lambda\to S(\lambda)$ est holomorphe au voisinage
% ms. line 14
de $0$, avec un rayon de convergence égal à $1$,
% ms. line 15
\deleted{\unsure{ce qui fait que ses|ce qui fait que sur}} \inserted{et ses restrictions} aux sous-espaces vectoriels
% ms. line 16
complexes de $\mathrm{dim}$ finie sont holomorphes.
% ms. line 17
Cependant, \deleted{elle} cette fonction n'est pas bornée
% ms. line 18
dans la boule unité de $l^1$, et n'est par
% ms. line 19
suite pas une fonction entière de $\lambda$.

% ms. line 20
\uline{Corollaire} On a \inserted{pour $\lambda\in l^1$} $\sqrt[n]{n!|\alpha_n(\lambda)|}\to 0$,
% Vertical rule beside the corollary.
% ms. line 21
$\sqrt[n]{|\alpha_n(\lambda)|}=o\left(\frac{1}{n}\right)$, $\alpha_n(\lambda)=\frac{\varepsilon_n^n}{n!}$ avec $\varepsilon_n\to 0$
% ms. line 22
Si $\sum a_nz^n$ est une fonction entière de
% ms. line 23
genre zéro, la suite des $a_n$ jouit
% ms. line 24
des $=^{\mathrm s}$ propriétés.

% ms. line 25
En effet, $\sqrt[n]{n!|\alpha_n(\lambda)|}\to 0$ exprime que
% archive sheet number: 121
